\section{A correlated singular realization with a continuous roof}
\label{sec:v59-markov-realization}
The finite-digit example of Section~\ref{sec:v58-pressure-contact}
does not test a continuous density or correlated symbols. We now verify
the same pressure mechanism on a two-state Markov baker transformation.
Its stable coordinate gives an actual continuous roof and, in the next
section, a pointwise density theorem. The comparison is with finite-state
Markov additive local limits \cite{HL}; the additional calculation here
keeps the continuous deterministic coboundary and its boundary sources
inside the coarea formula. It is not a claim of a new general Markov
local limit theorem.

Put
\begin{equation}\label{eq:v59-markov-data}
 P=\begin{pmatrix}3/4&1/4\\1/3&2/3\end{pmatrix},\qquad
 \pi=(4/7,3/7),\qquad \rho_0=3/4.
\end{equation}
The state space is the disjoint union of two unit squares, with
$\mu(\{i\}\times dx\,dy)=\pi_i dx\,dy$. Let
$a_{i0}=0$, $a_{i1}=P_{i0}$, $b_{j0}=0$, and $b_{j1}=P_{j0}$.
On $a_{ij}\le x<a_{ij}+P_{ij}$ define
\begin{equation}\label{eq:v59-markov-baker}
 F(i,x,y)=\left(j,\frac{x-a_{ij}}{P_{ij}},b_{ji}+P_{ji}y\right).
\end{equation}
Values on partition boundaries are immaterial. Detailed balance
$\pi_iP_{ij}=\pi_jP_{ji}$ holds. Consequently the images partition each
target square horizontally and the Jacobian $P_{ji}/P_{ij}$ preserves
$\mu$. The map is invertible modulo boundaries, expands the first
coordinate by at least $4/3$ and contracts the second by at most $3/4$.
Its cuts are genuine vertical and horizontal singularity lines.
The state sequence $I_r$ has transition matrix $P$. In particular
\begin{equation}\label{eq:v59-nonindependence}
 \operatorname{Cov}_\mu(I_0,I_r)=\frac{12}{49}(5/12)^r.
\end{equation}
Thus these symbols are not independent. This statement concerns the
specified coding, not absence of a Bernoulli isomorphism.

For a small real $\epsilon$ set
\begin{align}
 \xi_\epsilon(i,j)&=j+\epsilon ij,\qquad
 \bar\xi_\epsilon=3/7+2\epsilon/7,\notag\\
 \tau_\epsilon(i,x,y)&=3+\xi_\epsilon(i,j)+y(F(i,x,y))-y.
                                             \label{eq:v59-continuous-roof}
\end{align}
For $|\epsilon|<1/4$ this is a bounded positive roof, piecewise affine
in $y$ with a nonzero derivative. No independent smoothing variable is
added. Along the original deterministic trajectory,
\begin{equation}\label{eq:v59-roof-telescoping}
 S_m\tau_\epsilon=m(3+\bar\xi_\epsilon)
 +\sum_{r=0}^{m-1}(\xi_\epsilon(I_r,I_{r+1})-\bar\xi_\epsilon)
 +y_m-y_0.
\end{equation}

\subsection{An infinite-dimensional physical operator}
Let $\mathcal B$ be pairs of complex Lipschitz functions on $[0,1]$,
with norm $\|h\|_\infty+\max_i\operatorname{Lip}(h_i)$. Define
\begin{align}
 (\mathcal P_{\epsilon,z}h)_i(y)
 &=\sum_j P_{ij}e^{iz(\xi_\epsilon(i,j)-\bar\xi_\epsilon)}
                       h_j(b_{ji}+P_{ji}y),\notag\\
 \mathcal Q_{\epsilon,z}
 &=M_{e^{-izy}}\mathcal P_{\epsilon,z}M_{e^{izy}}.
                                             \label{eq:v59-markov-operator}
\end{align}
Conditional integration over the future $x$ coordinate proves the
exact scalar identity, for bounded initial tests $a(i,y)$ and
Lipschitz terminal tests $d(i,y)$,
\begin{equation}\label{eq:v59-physical-pairing}
 \int a e^{iz(S_m\tau_\epsilon-m(3+\bar\xi_\epsilon))}
                 (d\circ F^m)d\mu
       =\sum_i\pi_i\int_0^1a_i(y)(\mathcal Q_{\epsilon,z}^md)_i(y)\,dy.
\end{equation}
This is a physical conditional-expectation operator for the stated
endpoint class, not a claim that it is the complete billiard operator.

\begin{lemma}[Constants, contraction and the visible eigenvalue]
\label{lem:v59-markov-spectrum}
Near $(\epsilon,z)=(0,0)$ and $(0,2\pi)$, the family
$\mathcal Q_{\epsilon,z}$ is strongly analytic in $z$ and has a simple
leading eigenvalue, an analytic rank-one projection, and a complementary
power bound $C\rho_1^m$ with a common $\rho_1<1$. The leading eigenvalue
is the one from the matrix
\begin{equation}\label{eq:v59-markov-matrix}
 D_\epsilon(z)_{ij}
       =P_{ij}e^{iz(\xi_\epsilon(i,j)-\bar\xi_\epsilon)}.
\end{equation}
All constants are uniform for $\epsilon$ in a sufficiently small closed
interval. The continuous roof has physical witnesses of nonzero
amplitude near its resonance.
\end{lemma}
\begin{proof}
On compact complex sets the exponential and its derivatives are bounded
multipliers of $\mathcal B$, proving strong analyticity. The subspace
of functions constant in $y$ is invariant for $\mathcal P$, with matrix
$D_\epsilon(z)$. On its quotient the Lipschitz seminorm is a norm, and
\[
 \operatorname{Lip}(\mathcal P_{\epsilon,z}^mh)
 \le(\rho_0e^{C|\operatorname{Im}z|})^m\operatorname{Lip}(h).
\]
Choose the complex neighborhoods so that the base is less than one.
Splitting a function into its values at $y=0$ and a function vanishing
there writes the operator in upper triangular blocks. Its quotient
block has the displayed spectral-radius bound. Outside that disk its
spectrum is exactly that of $D_\epsilon(z)$, by the block resolvent
formula. At $(0,0)$ the two matrix eigenvalues are $1,5/12$; at
$(0,2\pi)$ they are these numbers multiplied by $e^{-6\pi i/7}$.
Small fixed resolvent contours and the Neumann series therefore yield
the stated splitting and uniform complementary powers. Conjugation by
$e^{izy}$ preserves them on the fixed neighborhoods.

At $\epsilon=0,z=2\pi$ choose the physical tests
$a(i,y)=e^{2\pi iy}$ and $d(i,y)=e^{-2\pi iy}$. The two endpoint
factors cancel $y_m-y_0$ in \eqref{eq:v59-roof-telescoping}; the
pairing is exactly $e^{-6\pi im/7}$. Its leading amplitude is one.
Analyticity of the scalar pairing and projection keeps this amplitude
bounded away from zero on a smaller neighborhood. These tests are
fixed, not changed with the complex tilt. The remainder comes from
the actual operator splitting just proved.
\end{proof}

\subsection{Pressure and a genuinely damped peak}
Let $\Lambda(s,t)$ be the Perron eigenvalue near one of
\[
 \begin{pmatrix}3/4&(1/4)e^s\\1/3&(2/3)e^{s+t}\end{pmatrix},
 \qquad C(s,t)=\log\Lambda(s,t).
\]
Differentiating its characteristic polynomial
\begin{equation}\label{eq:v59-characteristic-polynomial}
 \Lambda^2-(3/4+(2/3)e^{s+t})\Lambda
                     +(1/2)e^{s+t}-(1/12)e^s=0
\end{equation}
gives
\begin{equation}\label{eq:v59-joint-covariance}
 DC(0,0)=(3/7,2/7),\qquad
 D^2C(0,0)=\frac1{343}
       \begin{pmatrix}204&192\\192&198\end{pmatrix}.
\end{equation}
Its determinant is $72/2401>0$. In particular the centered observable
$\xi_\epsilon$ has asymptotic variance
\[
 \sigma_\epsilon^2=(204+384\epsilon+198\epsilon^2)/343.
\]
The same variance belongs to the continuous roof, since the difference
of the two centered sums in \eqref{eq:v59-roof-telescoping} is bounded
and the symbolic variance is $O(m)$.

\begin{theorem}[Pressure contact for the correlated continuous roof]
\label{thm:v59-markov-contact}
The continuous roof in \eqref{eq:v59-continuous-roof} satisfies all the
physical pressure, visibility and contact hypotheses of
Theorem~\ref{thm:v58-flat-contact} near the peak at $2\pi$. Its real
maximizer, damping and centered drift satisfy
\begin{align}
 a_\epsilon&=2\pi-\frac{32\pi}{17}\epsilon+O(\epsilon^2),\notag\\
 \kappa_\epsilon&=\frac{12\pi^2}{119}\epsilon^2+O(\epsilon^3),\notag\\
 v_\epsilon&=-\frac{3456\pi^2}{99127}\epsilon^2+O(\epsilon^3).
                                              \label{eq:v59-markov-expansions}
\end{align}
Here $D_z\psi_\epsilon(a_\epsilon)=iv_\epsilon$ for the logarithm
of the \emph{centered} pairing. Thus $v_\epsilon^2=o(\kappa_\epsilon)$,
and the peak has positive damping for all sufficiently small nonzero
$\epsilon$. This realization includes a continuous deterministic roof;
it is not an independent finite-digit exponential sum.
\end{theorem}
\begin{proof}
For real $t$ the matrix $D_\epsilon(it)$ is positive. Its simple Perron
eigenvalue gives a real analytic pressure $\mathscr P_\epsilon(t)$
with zero value and gradient at zero and Hessian $\sigma_\epsilon^2$.
Uniform Perron bounds give the exponential-moment upper estimate for
the symbolic sum. The bounded factor $e^{-t(y_m-y_0)}$ in the original
roof changes only its prefactor by $e^{|t|}$, not its pressure. This
checks the pressure on the actual probability. Visibility was proved in
Lemma~\ref{lem:v59-markov-spectrum}.

At $\epsilon=0$ the identity
$D_0(2\pi+w)=e^{-6\pi i/7}D_0(w)$ holds for all complex $w$.
Consequently the Hessian at resonance is exactly the physical variance.
The real logarithm has a strict quadratic maximum there, so the implicit
function theorem gives the real maximizer $a_\epsilon$.

Here are the coefficient checks without a numerical fit. Put
$z=2\pi+q$. Up to the fixed logarithm branch, the centered logarithm is
\[
 \psi_\epsilon(z)=-iz(3/7+2\epsilon/7)
                                    +C(iq,iz\epsilon).
\]
The linear terms reduce to $-6\pi i/7$. Minimizing the quadratic
form in \eqref{eq:v59-joint-covariance} gives
$q=-2\pi(192/204)\epsilon+O(\epsilon^2)$ and the residual variance
$\Sigma_{22}-\Sigma_{12}^2/\Sigma_{11}=6/119$.
This proves the first two expansions. The three derivatives needed for
the drift are
\[
 C_{sss}(0)=4908/16807,\quad C_{sst}(0)=9600/16807,
                       \quad C_{stt}(0)=14016/16807.
\]
Inserting $w=(-32\pi/17,2\pi)$ into the derivative of the cubic term
$-iD^3C(0)[(q,z\epsilon)]^3/6$ gives
$-D^3C(0)[(1,0),w,w]/2=-3456\pi^2/99127$.
The quadratic term contributes only to the real derivative, which is
zero at the real maximum. The remaining terms are $O(\epsilon^3)$.
These derivatives follow directly from
\eqref{eq:v59-characteristic-polynomial}. Uniform analyticity justifies
the remainders. After reducing the parameter interval, the positive
quadratic damping shows that its only zero is $\epsilon=0$, where exact
contact has been checked. All compact-support hypotheses now hold.
\end{proof}
